\documentclass[12pt]{article}
\usepackage[utf8]{inputenc}
\usepackage[T1]{fontenc}
\usepackage{hyperref}
\title{Toki Pona Language Resources for K-12 Curriculum Development}
\author{Compiled by Claude Sonnet 5 (claude-sonnet-5) for Andr\'as Kornai}
\date{}
\begin{document}
\maketitle

\section{Introduction}
Toki Pona is a constructed language created by Sonja Lang. It has no native-speaking community in the traditional sense, and its speakers are learners and hobbyists. This report is based on one search round. Almost all of the resources found come from the community-run kama sona wiki. No school curriculum or school programme was found.

\section{Primary Educational Resources}
\begin{itemize}
\item Recommended learning resources (kama sona wiki), a curated list of courses, flashcards and videos: \url{https://sona.pona.la/wiki/learn}
\item Learning resources (kama sona wiki): \url{https://sona.pona.la/wiki/resources}
\item \emph{Toki Pona: The Language of Good} (lipu pu), Sonja Lang's lesson book with exercises, described on the wiki: \url{https://sona.pona.la/wiki/Toki_Pona:_The_Language_of_Good?oldid=2394}
\end{itemize}
The wiki also names \emph{lipu sona pona}, \emph{lipu sona}, the video series \emph{o pilin e toki pona}, and an Anki deck, but I did not see a direct URL for each of them. They are not listed as separate items.

\section{Linguistic and Reference Materials}
The wiki describes \emph{Toki Pona Dictionary} (lipu ku) by Sonja Lang, the dictionary sites nimi.pona.la and Linku, and the grammar \emph{nasin toki pona}. No separate URLs for these were seen in the search results beyond the wiki pages above.

\section{Religious and Community Resources}
The wiki points to the Reddit community r/tokipona and the kama sona Discord server. I did not see URLs for them.

\section{Digital and Audio Resources}
Digital resources are the online dictionaries and courses listed on the wiki pages above. A secondary overview article is at \url{https://blog.openl.io/toki-pona/}.

\section{Curriculum Development Context}
The language has an official lesson book and community-built courses, but no school curriculum was found. The wiki notes that an error-laden video series is not recommended, so materials should be checked against the wiki's current guidance. A K-12 curriculum would need to be built from the community materials and the official books.
\end{document}